% chapters/retrieval.tex — Retrieval chapter (narrative text only)
\chapter{Retrieval-Augmented Reasoning}
\label{chap:retrieval}

Retrieval-augmented generation (\RAG) grounds \LLM{} outputs in external
knowledge sources, reducing hallucination and enabling the agent to reason
over information beyond its training cut-off~\cite{wu2025agenticreasoningtools}.
\index{retrieval-augmented generation}
\index{RAG}

\section{Reactive vs.\ Proactive Retrieval}
\label{sec:retrieval-proactive}

Early \RAG systems retrieved documents only in response to an explicit
query at inference time.
Recent work on proactive retrieval~\cite{wang2025proactiveretrievalmedical}
anticipates information needs before the \LLM{} formulates a query,
pre-fetching relevant context and reducing latency on multi-step tasks.

\section{Retrieval in Medical Multimodal Settings}
\label{sec:retrieval-medical}

\cite{wang2025proactiveretrievalmedical} demonstrated that proactive
retrieval improves diagnostic accuracy when applied to multimodal medical
\LLM{}s that must reason over both clinical text and radiological images.
The key insight is that retrieval decisions can be made by a lightweight
routing model that monitors agent state, rather than waiting for the base
\LLM{} to request information explicitly.

\section{Agentic Retrieval Frameworks}
\label{sec:retrieval-agentic}

\cite{wu2025agenticreasoningtools} frame retrieval as an agentic sub-policy:
the agent learns when to retrieve, what to retrieve, and how to integrate
retrieved content into its current reasoning chain.
This formulation connects retrieval to the policy objective in
Equation~\eqref{eq:planning-objective} of Chapter~\ref{chap:planning}: the
agent maximises expected reward by selecting retrieval actions only when the
information gain justifies the latency cost.
